\section{Historical Metadata and Acquisition Feasibility}
\label{sec:temporal_authority}
The historical companion asks what retained registry metadata can support before an event cutoff. It is separate from the cross-sectional reverse-reach experiment. It does not test whether the primary rankings would have predicted historical compromise.

\subsection{Bounded Cutoff-SemVer Reconstruction}
For each retained root record, the reconstruction chooses the highest currently retained, non-prerelease version satisfying a runtime dependency range and published by the root cutoff. Expansion stops at depth three or 250 nodes per root. Peer, optional, development, bundled, workspace, override, git, and tarball dependency behavior is excluded. Deleted historical versions may be unavailable; present-day retention cannot recreate every past registry state.

Table~\ref{tab:q1_temporal_reconstruction} summarizes the bounded reconstruction.
% Generated by scripts/reconcile_full_manuscript.py
\begin{table*}[t]
\centering
\footnotesize
\caption{Retained bounded registry-SemVer reconstruction manifest. The recorded resolution share is conditional on attempted expansions within the bounds.}
\label{tab:q1_temporal_reconstruction}
\begin{tabular}{lr}
\toprule
Recorded dimension & Value\\
\midrule
Retained root records & 45\\
Distinct event/root/cutoff tuples & 44\\
Recorded edge attempts & 1,848\\
Resolved recorded edges & 1,841\\
Node counts summed over roots & 1,272\\
Unresolved expansions & 7\\
Resolved share of recorded attempts & 99.62\%\\
Maximum expansion depth & 3\\
Node cap per root & 250\\
Roots reaching cap & 1\\
\bottomrule
\end{tabular}
\par\smallskip\begin{minipage}{0.98\textwidth}\footnotesize Only regular dependencies are included; depth is capped at three and nodes at 250 per root. Highest retained cutoff-valid satisfying versions are used. These are bounded approximations, not lockfile reconstructions or completeness guarantees. Raw-response checksums establish retained-object identity, not publisher authenticity.\end{minipage}
\end{table*}


The source contains 45 root records but 44 distinct event/root/cutoff tuples: one \texttt{node-ipc} tuple is repeated. We report retained-record counts without treating them as independent incidents. Of 1,848 recorded edge attempts, 1,841 resolve and seven have unavailable target registries, yielding 99.62\% among recorded attempts. This percentage excludes unvisited expansions beyond the bounds. The 1,272 node count is summed across roots, not globally deduplicated. One record reaches the node cap. These measurements do not establish full-tree completeness, downstream population reconstruction, or unbiased historical rankings.

\subsection{Release-Field Availability and Authority Limits}
Table~\ref{tab:q1_release_authority} reports the availability of release fields.
% Generated by scripts/reconcile_full_manuscript.py
\begin{table*}[t]
\centering
\footnotesize
\caption{Field availability in 45 retained pre-event version-object records. Comparison availability is not continuity, equality, or authenticated release authority.}
\label{tab:q1_release_authority}
\begin{tabular}{lrr}
\toprule
Metadata availability & Count & Share\\
\midrule
Retained selected-version objects & 45 & 100.0\%\\
Immediate predecessor available & 44 & 97.8\%\\
Repository-field comparison available & 42 & 93.3\%\\
gitHead-field comparison available & 34 & 75.6\%\\
publishConfig comparison available & 11 & 24.4\%\\
Selected dist.integrity present & 45 & 100.0\%\\
Selected embedded maintainers present & 45 & 100.0\%\\
\bottomrule
\end{tabular}
\par\smallskip\begin{minipage}{0.98\textwidth}\footnotesize The 45 records include a repeated event/root/cutoff tuple; they are not independent incidents. Embedded maintainer metadata is not a signature. No version-specific provenance endpoint requests, bundle retrievals, signature verifications, or attestation-state classifications were performed. Historical attestation presence/absence is unknown, not 0/45 absent.\end{minipage}
\end{table*}


The availability audit finds predecessor metadata for 44 of 45 retained records, repository-field comparison availability for 42, and \texttt{gitHead}-field comparison availability for 34. Availability does not show that URLs stayed equal, commits were valid or available, or a publisher retained authority. The audit classified no field deltas, built no authority scores, and evaluated no release-gating control.

All 45 records contain \texttt{dist.integrity} and embedded maintainer metadata. A distribution hash identifies expected artifact bytes; an embedded maintainer array is not a cryptographic signature. Neither proves a benign or reproducible build. The attestation access assessment made no version-specific requests, retrieved no bundles, ran no signature verification, and classified no attestation states. Historical attestation presence is therefore \emph{unknown}; zero retrievals cannot be reported as zero attestations. Signing studies motivate the need for stronger evidence~\cite{okafor2025,schorlemmer2024}, without supplying the missing measurements for this cohort.

The LiteLLM source-tag/artifact discrepancy~\cite{litellmincident2026} illustrates why artifact identity, intended source, and release authority must be evaluated separately. Contextual reports do not add observations to the 45-record companion.

\subsection{Metadata-Acquisition Pilot}
A separate metadata pilot sampled 100 identifiers uniformly without replacement from the original 25,228-identifier frame (seed 20261020). Of 92 ordinary-shaped names, 90 yielded metadata and two returned HTTP 404; all eight bundled-context identifiers returned HTTP 405. No identifier was replaced. The acquisition cutoff was October 20, 2026, 00:48:38 Asia/Dhaka. The pilot characterizes retrieval feasibility and sampling-frame defects; it includes no measured incident outcomes or forecasting evaluation.
